\subsection{原子主导 residue 单纯形轨道与充分统计外置的覆盖坍缩}\label{subsec:conclusion-fullresidue-simplex-sufficient-externalization-collapse}
沿用定理 \ref{thm:xi-time-part65c-atomic-centered-residue-simplex}、推论 \ref{cor:xi-time-part65c-atomic-centered-residue-isotropic}、定理 \ref{thm:xi-time-part65c-atom-dominated-normalized-residue-profile}、推论 \ref{cor:xi-single-atom-centered-residue-l2-exponent-truncation} 与命题 \ref{prop:pom-sufficient-statistic-residual-noninvertibility} 的记号。前述链条已经把单原子 residue 扰动的 centered 几何、full/core 过渡下的精确指数与充分统计外置的不可能性分别固定下来。把这些局部结果合并后，还可进一步抽出如下终端后果：在原子主导区，偶长度归一化 residue-law 的偏差轨道并非一般零均值云，而是落在一个边长与模数无关的正则单纯形上；与此同时，任何只依赖充分统计量的外置机制在长纤维上的最优可逆覆盖率都会指数坍缩。

\begin{theorem}[原子主导区中归一化 residue-law 的正则单纯形轨道]\label{thm:conclusion-full-residue-regular-simplex-orbit}
\leanverified{paper\_conclusion\_full\_residue\_regular\_simplex\_orbit}
固定模数 \(m\ge 2\)，并假设
\[
\rho_m^{\mathrm{core}}<1.
\]
定义偶长度归一化 residue-law 偏差
\[
\nu_q
:=
\frac{a_{2q}^{\mathrm{full}}(1)}{2}\bigl(\mu_{m,\bullet}^{\mathrm{full}}(2q)-u_m\bigr)
\in \RR^{\ZZ/m\ZZ},
\qquad
u_m:=\frac1m\mathbf 1_m.
\]
对每个 \(r\in \ZZ/m\ZZ\) 再记
\[
v_r:=e_r-u_m.
\]
则有极限轨道公式
\[
\nu_q=v_{q\bmod m}+o(1).
\]
从而极限支撑恰为
\[
\mathcal V_m:=\{v_r:\ r\in\ZZ/m\ZZ\},
\]
并且 \(\mathcal V_m\) 在零和超平面
\[
H_m:=\left\{x\in\RR^{\ZZ/m\ZZ}:\ \sum_{r\in\ZZ/m\ZZ}x_r=0\right\}
\]
中构成一个正则 \(A_{m-1}\)-单纯形，满足
\[
\langle v_r,v_s\rangle=\delta_{rs}-\frac1m,
\qquad
\|v_r\|_2^2=\frac{m-1}{m},
\qquad
\|v_r-v_s\|_2^2=2
\quad (r\neq s).
\]
特别地，其边长恒等于 \(\sqrt2\)，与模数 \(m\) 无关。

进一步，偶子列的经验测度收敛到 \(\mathcal V_m\) 上的均匀原子测度：
\[
\frac1N\sum_{q=0}^{N-1}\delta_{\nu_q}
\Longrightarrow
\frac1m\sum_{r\in\ZZ/m\ZZ}\delta_{v_r}.
\]
相应二阶矩阵满足
\[
\frac1N\sum_{q=0}^{N-1}\nu_q\nu_q^{\mathsf T}
\longrightarrow
\frac1m I_m-\frac1{m^2}J_m,
\]
其中 \(J_m\) 为全 \(1\) 矩阵；因此在 \(H_m\) 上，该极限二阶矩恰为
\[
\frac1m\,\mathrm{Id}_{H_m}.
\]
\end{theorem}
\begin{proof}
由定理 \ref{thm:xi-time-part65c-atom-dominated-normalized-residue-profile}，
\[
\frac{a_{2q}^{\mathrm{full}}(1)}{2}\bigl(\mu_{m,\bullet}^{\mathrm{full}}(2q)-u_m\bigr)
=
e_{r_q}-u_m+o(1),
\qquad
r_q\equiv q\pmod m.
\]
这正是
\[
\nu_q=v_{q\bmod m}+o(1).
\]

再由定义直接计算，
\[
\langle v_r,v_s\rangle
=
\langle e_r-u_m,e_s-u_m\rangle
=
\delta_{rs}-\frac1m,
\]
故
\[
\|v_r\|_2^2=1-\frac1m=\frac{m-1}{m}.
\]
若 \(r\neq s\)，则
\[
v_r-v_s=e_r-e_s,
\]
因而
\[
\|v_r-v_s\|_2^2=\|e_r-e_s\|_2^2=2.
\]
于是 \(\mathcal V_m\) 为嵌入 \(H_m\) 的正则单纯形。

由于 \(q\mapsto q\bmod m\) 在区间 \(\{0,\dots,mN-1\}\) 上均匀遍历 \(\ZZ/m\ZZ\)，故经验测度极限为
\[
\frac1m\sum_{r\in\ZZ/m\ZZ}\delta_{v_r}.
\]
进而
\[
\frac1m\sum_{r\in\ZZ/m\ZZ}v_rv_r^{\mathsf T}
=
\frac1m\sum_{r}(e_r-u_m)(e_r-u_m)^{\mathsf T}
=
\frac1m I_m-\frac1{m^2}J_m.
\]
而 \(J_m\) 在零和超平面 \(H_m\) 上作用为零，故其限制为 \(\frac1m\mathrm{Id}_{H_m}\)。证毕。
\end{proof}

\begin{corollary}[原子主导 residue-law 的统一指数冻结]\label{cor:conclusion-full-residue-phi-square-universal-exponent}
在定理 \ref{thm:conclusion-full-residue-regular-simplex-orbit} 的设定下，偶长度与奇长度偏差分别满足
\[
\lim_{q\to\infty}
\frac{a_{2q}^{\mathrm{full}}(1)}{2}
\bigl\|\mu_{m,\bullet}^{\mathrm{full}}(2q)-u_m\bigr\|_2
=
\sqrt{\frac{m-1}{m}},
\]
以及
\[
\lim_{q\to\infty}
a_{2q+1}^{\mathrm{full}}(1)
\bigl\|\mu_{m,\bullet}^{\mathrm{full}}(2q+1)-u_m\bigr\|_2
=
0.
\]
更进一步，指数衰减率完全脱离模数 \(m\)：
\[
\limsup_{n\to\infty}
\bigl\|\mu_{m,\bullet}^{\mathrm{full}}(n)-u_m\bigr\|_2^{1/n}
=
\varphi^{-2}.
\]
因此，一旦进入 \(\rho_m^{\mathrm{core}}<1\) 的原子主导区，residue-law 的混合指数就在全部模数上冻结为同一个常数 \(\varphi^{-2}\)；模数依赖仅保留在前导幅值的正则单纯形几何中。
\end{corollary}
\begin{proof}
第一式由定理 \ref{thm:conclusion-full-residue-regular-simplex-orbit} 与
\[
\|e_r-u_m\|_2=\sqrt{\frac{m-1}{m}}
\]
直接得到。奇长度结论则正是定理 \ref{thm:xi-time-part65c-atom-dominated-normalized-residue-profile} 的奇长度部分。

最后，由推论 \ref{cor:xi-single-atom-centered-residue-l2-exponent-truncation}，
\[
\limsup_{n\to\infty}
\bigl\|\mu_{m,\bullet}^{\mathrm{full}}(n)-u_m\bigr\|_2^{1/n}
=
\frac{\max\{1,\rho_m^{\mathrm{core}}\}}{\lambda}.
\]
在 \(\rho_m^{\mathrm{core}}<1\) 下，上式化为
\[
\lambda^{-1}=\varphi^{-2},
\]
其中 \(\lambda=\rho(M(1))=\varphi^2\)。证毕。
\end{proof}

\begin{theorem}[充分统计量外置的最优可逆覆盖率坍缩]\label{thm:conclusion-sufficient-statistic-externalization-success-collapse}
\leanverified{paper\_conclusion\_sufficient\_statistic\_externalization\_success\_collapse}
固定 \(m\ge 1\) 与纤维 \(x\in X_m\)。设
\[
\Gamma_x\simeq\bigsqcup_i P_{\ell_i}
\]
为该纤维的规范分量分解，并在纤维 \(\Fold_m^{-1}(x)\) 上考虑任意仅依赖充分统计量 \(r(\omega)\) 的外置残差
\[
\mathcal R(\omega)=g(r(\omega)).
\]
若在该纤维上赋予均匀先验，则任意确定性译码器基于 \((x,\mathcal R(\omega))\) 的最优成功概率满足
\[
p_{\mathrm{succ}}(x)
\le
\frac{d_x+1}{d_m(x)}
=
\frac{1+\sum_i\lceil \ell_i/2\rceil}{\prod_i F_{\ell_i+2}}.
\]
特别地，在单分量情形 \(\Gamma_x\simeq P_\ell\) 下，
\[
p_{\mathrm{succ}}(x)
\le
\frac{\lceil \ell/2\rceil+1}{F_{\ell+2}}
=
O(\ell\,\varphi^{-\ell}).
\]
因此，充分统计量外置不只是不可能在一般纤维上全可逆；在长的不可约纤维上，其最优可逆覆盖率本身就以指数速度坍缩。
\end{theorem}
\begin{proof}
由命题 \ref{prop:pom-sufficient-statistic-residual-noninvertibility}，
\[
\bigl|\mathcal R(\Fold_m^{-1}(x))\bigr|\le d_x+1,
\qquad
d_m(x)=\prod_i F_{\ell_i+2},
\qquad
d_x=\sum_i\left\lceil\frac{\ell_i}{2}\right\rceil.
\]
在均匀先验下，任意确定性译码器对每个残差值至多能正确恢复一个原像，因此成功概率至多等于“可区分残差值数 / 纤维总大小”，即
\[
p_{\mathrm{succ}}(x)\le \frac{d_x+1}{d_m(x)}.
\]
代入上面的闭式便得到第一条不等式。

若 \(\Gamma_x\simeq P_\ell\)，则
\[
d_m(x)=F_{\ell+2},
\qquad
d_x=\left\lceil\frac{\ell}{2}\right\rceil,
\]
从而
\[
p_{\mathrm{succ}}(x)
\le
\frac{\lceil \ell/2\rceil+1}{F_{\ell+2}}.
\]
再用 Fibonacci 增长
\[
F_{\ell+2}\asymp \varphi^\ell
\]
便得
\[
p_{\mathrm{succ}}(x)=O(\ell\,\varphi^{-\ell}).
\]
证毕。
\end{proof}

\begin{corollary}[完备充分统计外置的全局可逆覆盖比例坍缩]\label{cor:conclusion-sufficient-statistic-global-coverage-collapse}
\leanverified{paper\_conclusion\_sufficient\_statistic\_global\_coverage\_collapse}
固定分辨率 \(m\ge 1\)。设外置机制在每条纤维上仅依赖逆代换次数 \(r(\omega)\)，并定义
\[
\Psi:\Omega_m\to X_m\times\{0,1\}^\ast,
\qquad
\Psi(\omega):=\bigl(\Fold_m(\omega),\mathcal R(\omega)\bigr).
\]
则其可逆覆盖比例满足
\[
\frac{|\Psi(\Omega_m)|}{2^m}
\le
\frac{(m+1)|X_m|}{2^m}.
\]
由于
\[
|X_m|=F_{m+2}=\Theta(\varphi^m),
\]
遂有
\[
\frac{|\Psi(\Omega_m)|}{2^m}
=
O\!\left(m\left(\frac{\varphi}{2}\right)^m\right)\xrightarrow[m\to\infty]{}0.
\]
因此，只要外置坚持完全通过纤维后验的完备充分统计量进行压缩，则全局可逆覆盖比例必以指数速度塌缩到零。
\end{corollary}
\begin{proof}
对每个 \(x\in X_m\)，由定理 \ref{thm:conclusion-sufficient-statistic-externalization-success-collapse}，
\[
|\mathcal R(\Fold_m^{-1}(x))|\le d_x+1,
\qquad
d_x:=\deg Z_x.
\]
而由 \(r(\omega)\in\{0,1,\dots,m\}\)，可知 \(d_x\le m\)，故
\[
|\mathcal R(\Fold_m^{-1}(x))|\le m+1.
\]
于是
\[
|\Psi(\Omega_m)|
\le
\sum_{x\in X_m}|\mathcal R(\Fold_m^{-1}(x))|
\le
(m+1)|X_m|.
\]
两边除以 \(2^m\)，再代入 \(|X_m|=F_{m+2}=\Theta(\varphi^m)\)，即得
\[
\frac{|\Psi(\Omega_m)|}{2^m}
=
O\!\left(m\left(\frac{\varphi}{2}\right)^m\right)\to 0.
\]
证毕。
\end{proof}

\noindent
由此，single-atom residue 通道与 sufficient-statistic 外置虽然都只涉及纤维内部的有限数据，却导出两种完全不同的刚性：前者在原子主导区把偶长度偏差压到一个与模数无关边长的正则单纯形轨道上，并在二阶矩意义下自动各向同性；后者则表明，哪怕只要求最优确定性回放，任何仅依赖充分统计量的外置账本在长纤维上也会出现指数级的可逆覆盖坍塌。于是，局部 residue 几何的高度规则性与充分统计外置的全局不可逆性，在本文框架中并不矛盾，而是同一纤维机制的两面。

\endinput
